%=====================================================================
% 第 5 节  另一类面算符
% 原文：Section 5. Another type of surface operators
%       （arXiv:1702.06499v2）
%=====================================================================
\section{另一类面算符}

\begin{keytakeaways}
\begin{itemize}
  \item \textbf{物理图像：}除余二维面算符之外，六维 $(2,0)$ 理论还可被余四维缺陷修饰。弦论中把 M5-M5'-膜系统沿 $x_7$ 方向拉开，会在其间接上若干 M2-膜，由此得到一个红外等价的 M5-M2 构型；M2 的端点扫出 $C$ 上的除子 $\sum_k P_k$。于是“在 $P$ 处插入一个余四维面算符”就对应于“在 $P$ 处插入一个退化场”。
  \item \textbf{数学工具：}Hanany--Witten 膜产生效应、AGT 对应在（推广的）共形块层面的翻译、仿射李代数 $\widehat{\mathfrak{sl}}_{2,k}$ 的模 $W_{-(k+h^{\vee})\la}$ 与 Virasoro 退化表示 $V_{-1/2b}$、Sklyanin 分离变量（SOV）方法，以及二维单值表示。
  \item \textbf{关键结论：}本节的核心猜想（式 \rf{Heckeconj}）是：在 $P\in C$ 处添加余四维面算符，等价于在共形块中 $P$ 处插入 $\widehat{\mathfrak{sl}}_{2,k}$-模 $W_{1/2\hbar}$。这正好把 Beilinson--Drinfeld 用 $W_{-(k+h^{\vee})\la}$ 描述的 Hecke 函子翻译成物理语言，$\la=1/2$ 对应 ${}^LG=\mathrm{PSL}(2,\C)$ 的基本表示。
  \item \textbf{证据与缺口：}证据来自 M5-M5' 与 M5-M2 的红外对偶、$W_{1/2\hbar}$ 在临界能级极限下的因子化、以及它产生二维单值表示这一唯一性论证；关键的局域化检验尚未进行。
\end{itemize}
\end{keytakeaways}

\bigskip
\noindent\rule{\textwidth}{0.6pt}
\bigskip

几何 Langlands 对应的关键要素之一是 Hecke 本征值性质。为了给下一节关于 Hecke 本征值性质的讨论作准备，我们在此讨论另一类面算符，它为我们提供了 Hecke 函子（Hecke functor）的物理实现。

\subsection{弦论动机}

从六维 $(2,0)$ 理论的角度看，研究用余四维缺陷修改该理论会产生什么效应，似乎是很自然的。这类缺陷构成对六维 $(2,0)$ 理论及其维数约化的自然探针，相关研究在 AGT 对应的语境中已经相当深入 \cite{AGGTV}。下面我们简要讨论：若同时考虑余二维与余四维的面算符，图像将如何改变。这类系统的存在性由弦论提示。特别地，依照 \cite{FGT} 的讨论，可以考虑 M 理论中支撑如下表所示的 M5-M5'-M2-膜构型：
\begin{center}
\begin{tabular}{l || c|c|c|c|c|c|c|c|c|c|c}
膜 & 0 & 1 & 2 & 3 & 4 & 5 & 6 & 7 & 8 & 9 & 10 \\ \hline\hline
M5  & x & x & x & x &   &   & x &   &   &   & x  \\
M5' & x & x &   &   &   &   & x &   & x & x & x \\
M2  & x & x &   &   &   &   &  &  x &   &   &   \\
\end{tabular}
\end{center}
具有这类支撑意味着：$\CM^4\times C$ 上的余四维面算符（其中 $\CM^4$ 局部地是一个如下形式的纤维积\footnote{记号 
$S_i^1$ 中的下标指的是用来定义参量为 $\ep_i$ 的 $\Om$-形变的那一个旋转对称性。} $\BR\times I\times
S^1_1\times S^1_2$）将表现为
\begin{itemize}
\item
$\CM^4$ 上支撑在 $\CM^2\simeq \BR\times S_{1}^1$ 上的面算符，
\item
$\BR\times I\times C$ 上 $N=4$ SYM 中的线算符，支撑在 $\BR\times \{\pi\}\times \{P\}$，
\item
或 $\BR\times C$ 上复 {\crd CS}（Chern--Simons）理论中的线算符，支撑在 $\BR\times \{P\}$。
\end{itemize}
由这些观察很容易看出，这类余四维面算符的存在会修改描述约化到 $\BR\times I$ 的二维 sigma 模型中的边界条件。下面我们简要描述：在作为二维 CFT 的有效描述中，余四维面算符的存在会导致哪些修改。

\subsection{对（推广的）AGT 对应带来的修改}

\cite{AGGTV} 论证：六维 $(2,0)$ 理论中余四维面算符的存在，在 $C$ 上紧化之后被描述为这样一个系统——一个 $4d$ 的 $\CS$ 类理论与 $\CM^4$ 中某个二维曲面 $\CM^2$ 上的二维规范线性 sigma 模型（GLSM）相耦合。相关 GLSM 的特征在于：其红外描述是一个以 Grassmann 流形为目标空间的非线性 sigma 模型。在最简单的 $\mathfrak{g}= \mathfrak{sl}_2$ 情形，例如得到的就是流向 $\mathbb{CP}^1$ sigma 模型的 GLSM。

关于在 CFT 内部、在余二维面算符存在时如何实现余四维面算符，我们的主要猜想如下：
\begin{equation}\label{Heckeconj}
\boxed{ 
\quad
\begin{aligned}
&\text{在 $P\in C$ 处添加一个余四维面算符}\\
& \text{由在 $P$ 处插入 $\widehat{\mathfrak{sl}}_{2,k}$-模 $W_{1/2\hbar}$ 来表示。}
\end{aligned}
\quad
}
\end{equation}
为了列举这一猜想的现有证据，先注意到它是由 \cite{FGT} 所述图像的一个轻微推广所提示的。该文献观察到，由 M5-M5'-膜组成的系统呈现出一种有趣的红外对偶，把它与一个 M5-M2-膜系统联系起来，其中若干 M2-膜位于 $C$ 上的各个点。这一红外对偶来自 Hanany--Witten 膜产生效应（brane creation effect）的一个变体，后者描述了一族相互红外等价的、通过把上述构型中已有的 M5-M5'-膜沿 $x_7$ 轴分开而得到的有效理论。为了得到红外等价的有效描述，需要引入额外的、悬挂在正在被分开的膜之间并沿 $x_7$ 轴延伸的 $M2$-膜。这些 $M2$-膜的端点定义了 $C$ 上的一个除子 $\sum_{k}P_k$。\cite{FGT} 论证：这一红外对偶表现为超对称配分函数之间的关系；借助 AGT 对应的各种推广，这些关系可翻译为仿射代数 $\widehat{\mathfrak{sl}}_{2,k}$ 与 $\mathfrak{Vir}_c$ 的共形块之间已知的关系。$\mathfrak{Vir}_c$ 的相关共形块是由在点 $P_1,\dots,P_h$（$h=4g-4+n$）处插入额外的、满足 $b^2=\hbar$ 的退化表示 $V_{-1/2b}$ 来定义的。这些关系来自 Sklyanin 分离变量（SOV）方法的一个推广。

\cite{FGT} 所考虑框架的一个简单变体是：在背景中再放入一个额外的“原初”（primordial）M2-膜，把它定域在 $P_0\in C$ 处并沿 $x_7$ 轴延伸，以此修改初始的 M5-M5'-膜系统。此时在红外对偶描述中，除 $P_1,\dots,P_h$ 处的插入之外，还会在 $P_0$ 处得到额外插入的一个退化表示 $V_{-1/2b}$。如 \cite{T10} 所解释的，这无非就是当在 $P$ 处额外插入 $\widehat{\mathfrak{sl}}_{2,k}$-模 $W_{1/2\hbar}$ 时、由 SOV 得到的 $\widehat{\mathfrak{sl}}_{2,k}$ 共形块的那个表示。

\subsection{更多证据}

另一种论证可以通过对 \cite{AGGTV} 中论证的适当推广得到：那里用退化场的插入来描述额外 M2-膜的作用。加入这样一个顶点算符，会使瞬子配分函数的空间再张量上 $\mathbb{CP}^1$ sigma 模型二维的真空向量空间。余四维面算符位置的变动，会在这个二维向量空间上生成一个单值表示。取 $\ep_2\ra 0 $ 的极限，会导致面算符上的自由度与体（bulk）中的自由度发生部分退耦，从而使瞬子配分函数分解为一个四维的奇异部分与一个二维的正则部分。$\widehat{\mathfrak{sl}}_{2,k}$-模 $W_{1/2\hbar}$ 是仿射李代数共形块中满足下述两条性质的唯一候选插入：(i) 生成一个二维单值表示；(ii) 在临界能级极限下，$W_{1/2\hbar}$ 的插入会从共形块中因子化出去。下一节我们将更详细地描述这一论证路线的若干要素。

若能通过局域化计算来检验我们的猜想 \rf{Heckeconj}，那将是很好的；这需要在两方面作推广：一方面是上面所引的关于余二维面算符的结果，另一方面是关于余四维面算符的已知结果（见 \cite{GLPP} 及其中的文献）。

在约化到 $N=4$ SYM 时，人们会用 $N=4$ SYM 中的 Wilson 与 't Hooft 线算符来表示余四维面算符 \cite{Wi12}。这些线算符与几何 Langlands 对应的 Hecke 函子之间的关系，已在 \cite{KW} 中作了广泛讨论。在 Cordova 与 Jafferis 于 \cite{CJ} 中所研究的约化方案中，将得到复 Chern--Simons 理论的 Wilson 线，这应能为 \rf{Heckeconj} 中所表述的猜想提供另一种视角。

最后可以指出，Beilinson 与 Drinfeld 工作中的 Hecke 函子，是由把 $\mathfrak{g}_k$ 的模 $W_{-(k+h^{\vee})\la}$ 插入共形块时所取的临界能级极限来描述的，其中 $\la$ 是 ${}^LG$ 的一个有限维表示的权，综述见 \cite[Section 20.5]{FBZ}。\rf{Heckeconj} 中出现的那些模则对应于 $\la=1/2$ 的情形，即 ${}^LG=\mathrm{PSL}(2,\C)$ 的基本表示。

\bigskip
\noindent\rule{\textwidth}{0.6pt}

% [Tutor Note: 式 \rf{Heckeconj} 的模与结尾 Beilinson--Drinfeld 公式的对应必须逐项核对：由 $k=-h^{\vee}-\ep_2/\ep_1$ 得 $k+h^{\vee}=-\ep_2/\ep_1=-1/\hbar$，即 $-(k+h^{\vee})=1/\hbar$，故 $\la=1/2$ 时 $-(k+h^{\vee})\la=\dfrac{1}{2\hbar}$，与 $W_{1/2\hbar}$ 相符——其中只用到本文的约定 $\hbar=\ep_1/\ep_2$。若换用文献中同样流行的 $\hbar=\ep_2/\ep_1$（并把能级写成 $k=-h^{\vee}-1/\hbar$），则同一个模会被写成 $W_{\hbar/2}$，于是 $1/(2\hbar)$ 与 $\hbar/2$ 这两套写法在两种约定下指的是同一个对象，直接照抄别人的公式时极易把指数写反。凡比较形如 $W_{\pm}$ 的下标，先固定 $\hbar$ 的定义再代换。]
% [Tutor Note: 退化表示 $V_{-1/2b}$ 的约定同样易错。取 Liouville/Virasoro 的标准记法 $\al_{m,n}=\frac{1-m}{2}b+\frac{1-n}{2}\frac{1}{b}$、$c=1+6(b+1/b)^2$，则 $V_{-1/2b}$ 对应 $\al_{1,2}=-1/(2b)$，而 $\al_{2,1}=-b/2$ 给出的是 $V_{-b/2}$；二者在 $b\ra1/b$（即 $\ep_1\leftrightarrow\ep_2$）下互换。正文给出 $b^2=\hbar$，故 $b\ra\infty$ 对应 $\ep_2\ra0$、$\hbar\ra\infty$、$k\ra-\infty$（临界能级），此时 $\al_{1,2}\ra0$，退化场趋于恒等场——这正是下面性质 (ii)“在临界能级极限下因子化”的另一种说法。]
% [Tutor Note: 表与正文的对应关系有两个易混点。其一，M5 与 M5' 在方向 7 上都没有支撑，因此沿 $x_7$ 把它们分开之后才可以悬挂 M2；M2 行给出的正是 0、1、7 三个方向，其中 0、1 是它在 $\CM^4$ 中留下的二维支撑，7 是被拉开的方向。其二，脚注中的 $S^1_1,S^1_2$ 是定义 $\Om$-形变（参数 $\ep_1,\ep_2$）的两个旋转对称，而第 4 节里余四维缺陷所缠绕的是 $S^1_\CS\subset\CM^4$ 与 $S^1_C\subset C$；把 $S^1_1$ 直接等同于 $S^1_\CS$ 缺乏依据，脚注的下标只在 $\CM^4$ 内部有意义。]
% [Tutor Note: “$W_{1/2\hbar}$ 是满足 (i)(ii) 的唯一候选”这一唯一性论断只是启发式的：性质 (i) 里的“二维单值表示”来自 $\mathbb{CP}^1$ GLSM 的两个真空，其单值群为 $\mathrm{GL}(2)$；性质 (ii) 则要求该模在临界能级极限下从共形块中分裂出来。二者合起来只说明“该模必须是二重退化且在临界能级退耦”，并不构成范畴层面的唯一性定理。真正的检验应如作者所言用局域化计算完成，而这一步在本文中并未进行，引用该猜想时不宜表述为已证。]
% [Tutor Note: 插入点数目 $h=4g-4+n$ 不要与 $\dim_{\C}\Bun_{SL_2}=3g-3+n$（即参数 $x$ 的个数）混为一谈；原文没有解释 $h$ 的来历，也没有给出 $\mathfrak{Vir}_c$ 的中心荷 $c$ 与 $k$、$\hbar$ 的关系式。读者若要把此处的共形块与第 3 节的 $\mathcal{CB}_{\rm te}(C,\hfg_k,\CE_x)$ 对齐，需要先补上 AGT/Liouville 字典中 $c=1+6(\sqrt{\hbar}+1/\sqrt{\hbar})^2$ 这一步，注意本文的 $b^2=\hbar$ 约定下 $c$ 随 $\hbar\ra\infty$ 发散，与 $\ep_2\ra0$ 的临界能级极限一致。]
